% ============================================================
\section{Methods}
\label{sec:methods}
% ============================================================

\subsection{Quantum-Inspired Feature Encodings}
\label{sec:encodings}

All three encodings are implemented as deterministic, parameter-free
transformations $\phi: \Rbb^d \to \Rbb^{\dout}$ fitted to training-set
statistics only.
No quantum simulator or circuit library is used; the encodings are pure NumPy
implementations of the geometric structure induced by the corresponding
quantum circuits.

\paragraph{Amplitude encoding.}
Given $\mathbf{x} \in \Rbb^d$, amplitude encoding produces
\begin{equation}
  \phi_{\mathrm{amp}}(\mathbf{x}) =
    \frac{\mathbf{x}}{\|\mathbf{x}\|_2 + \varepsilon}
    \;\text{(zero-padded to } 2^{\lceil\log_2 d\rceil}\text{)},
  \label{eq:amplitude}
\end{equation}
with $\varepsilon = 10^{-12}$ to guard against zero-norm inputs.
The output satisfies $\sum_i \phi_i^2 = 1$, mirroring the Born-rule
normalisation constraint on quantum state amplitudes.
The power-of-two padding expands dimensionality from $d$ to
$\dout = 2^{\lceil\log_2 d\rceil}$.

\paragraph{Angle encoding.}
Each feature $x_i$ is mapped to a pair of trigonometric values derived from
the quantum rotation gate
$R_y(\theta)|0\rangle = \cos(\theta/2)|0\rangle + \sin(\theta/2)|1\rangle$:
\begin{equation}
  \phi_{\mathrm{ang}}(\mathbf{x}) =
  [\cos(\theta_1/2),\,\sin(\theta_1/2),\,\ldots,\,
   \cos(\theta_d/2),\,\sin(\theta_d/2)],
  \quad \theta_i = \pi \tilde{x}_i,
  \label{eq:angle}
\end{equation}
where $\tilde{x}_i \in [-1,1]$ is obtained by min-max scaling each feature to
the training range.
The output dimension is $\dout = 2d$.
Each $(\cos,\sin)$ pair lies on the unit circle, satisfying
$\cos^2(\theta_i/2) + \sin^2(\theta_i/2) = 1$.

\paragraph{Basis encoding.}
Each continuous feature is quantised to $b = 8$ bits and unpacked to its
binary representation:
\begin{equation}
  \phi_{\mathrm{bas}}(\mathbf{x}) =
    [b_1^{(1)},\ldots,b_b^{(1)},\;
     b_1^{(2)},\ldots,b_b^{(2)},\;\ldots]
    \in \{0,1\}^{d \cdot b},
  \label{eq:basis}
\end{equation}
where $b_k^{(i)}$ is the $k$-th bit (MSB first) of the quantised value of
feature $x_i$.
The output dimension is $\dout = d \cdot b = 8d$.
This encoding is not differentiable; NTK analysis does not apply.

\subsection{Classical Baselines}
\label{sec:baselines}

Eight classical baselines are evaluated under the same matched-budget protocol:

\begin{itemize}
  \item \textbf{Raw linear:} standardised features with logistic regression.
  \item \textbf{RBF-SVM:} radial basis function kernel SVM with bandwidth
        tuned by cross-validation.
  \item \textbf{Polynomial (degree 2/3):} polynomial kernel logistic
        regression; degree chosen per dataset.
  \item \textbf{Random Fourier Features (RFF):} explicit kernel approximation
        of RBF with $\dout$ features~\citep{rahimi2007}.
  \item \textbf{PCA + logistic regression:} principal component projection
        to $\dout$ dimensions.
  \item \textbf{MLP (sklearn) and PyTorch MLP:} two-layer networks
        with ReLU activations; the strongest adaptive baselines.
\end{itemize}

All baselines receive the same output dimensionality $\dout$ as the QIE
encoding they are compared against, the same random seeds, and the same
train-test split.

\subsection{Datasets}
\label{sec:datasets}

Table~\ref{tab:datasets} summarises the ten benchmark datasets spanning six
categories.
Datasets were frozen before any experimental results were observed
(Phase~3 freeze).

\begin{table}[tbp]
\centering
\caption{Benchmark dataset roster.  Sample counts reflect the subset used in
experiments; full dataset sizes are given in parentheses where subsampling
was applied.}
\label{tab:datasets}
\small
\begin{tabular}{llrrrl}
\toprule
Dataset & Category & Samples & Features & Classes & Task \\
\midrule
UCI Wine           & Tabular     & 178                     & 13      & 3  & multiclass \\
UCI Breast Cancer  & Tabular     & 569                     & 30      & 2  & binary \\
Dry Bean           & Tabular     & 13{,}611                & 16      & 7  & multiclass \\
Credit Card Fraud  & Financial   & 284{,}807               & 30      & 2  & binary (imb.) \\
Fashion-MNIST$^*$ & Image       & 10{,}000 (of 70{,}000)  & 784     & 10 & multiclass \\
CIFAR-10$^*$      & Image       & 10{,}000 (of 60{,}000)  & 3{,}072 & 10 & multiclass \\
HIGGS              & Physics     & 500{,}000               & 28      & 2  & binary \\
High-dim parity    & Synthetic   & 10{,}000                & 20      & 2  & binary (XOR) \\
High-rank noise    & Synthetic   & 5{,}000                 & 200     & 2  & binary \\
Covertype$^\dagger$ & Large-scale & 50{,}000 (of 581{,}012) & 54    & 7  & multiclass \\
\bottomrule
\multicolumn{6}{l}{\small$^*$Stratified subsample; basis encoding at full scale expands to} \\
\multicolumn{6}{l}{\small\phantom{$^*$}6{,}272 (Fashion-MNIST) and 24{,}576 (CIFAR-10) dimensions,} \\
\multicolumn{6}{l}{\small\phantom{$^*$}making full-scale logistic regression computationally intractable.} \\
\multicolumn{6}{l}{\small$^\dagger$Stratified subsample for computational tractability.} \\
\end{tabular}
\end{table}

\subsection{Evaluation Protocol}
\label{sec:protocol}

All comparisons satisfy the following matched-budget constraints:
\begin{itemize}
  \item \textbf{Matched $\dout$:} classical baselines receive the same output
        dimensionality as the QIE encoding.
  \item \textbf{Matched optimisation:} identical solver, learning rate
        schedule, and epoch budget across all methods.
  \item \textbf{Matched regularisation:} same regularisation family and
        comparable grid search effort.
  \item \textbf{Matched seeds:} five independent seeds
        $\{7, 42, 99, 1337, 2026\}$ with identical train/test splits.
\end{itemize}

Primary performance metrics are accuracy and macro-averaged F1.
Statistical significance is assessed via paired $t$-test and Wilcoxon
signed-rank test over seeds, with effect size reported as Cohen's $d$.
The significance threshold is $\alpha = 0.05$.
Representation geometry is characterised by effective rank
$\erank(\mathbf{X}) = \exp\!\bigl(-\sum_i \hat{\sigma}_i \log \hat{\sigma}_i\bigr)$
(where $\hat{\sigma}$ are the normalised singular values), condition number
$\kappa$, and Centered Kernel Alignment (CKA)~\citep{kornblith2019}.
